% ============================================================================
%  GENERATED by scripts/gen-paper-tex.mjs from apps/web/lib/papers.ts
%  DO NOT EDIT BY HAND. Edit the manifest and regenerate:
%
%      node scripts/gen-paper-tex.mjs
%
%  The HTML at /papers/toronto-hardwood-climate-moisture-protocol and this document are the same manuscript.
%  That is the point: a hand-written .tex beside a manifest-driven page is two
%  documents that agree until the day they do not.
%
%  To publish the PDF:
%      1. compile this file (pdflatex, twice, for the table of contents)
%      2. mkdir -p apps/web/public/papers && mv <the-new>.pdf apps/web/public/papers/
%      3. git add apps/web/public/papers   -- in the same commit
%  The download button appears on its own; pdfIsPublished() checks the file.
% ============================================================================
\documentclass[11pt,a4paper]{article}
\usepackage[utf8]{inputenc}
\usepackage[T1]{fontenc}
\usepackage[margin=2.4cm]{geometry}
\usepackage{booktabs}
\usepackage{longtable}
\usepackage{array}
\usepackage{ragged2e}
\usepackage{enumitem}
\usepackage{xcolor}
\usepackage{mdframed}
\usepackage{titlesec}
\usepackage{fancyhdr}
\usepackage[hidelinks]{hyperref}

\definecolor{ecodark}{RGB}{16,20,28}
\definecolor{ecowood}{RGB}{145,95,45}
\definecolor{ecoaccent}{RGB}{0,100,120}
\definecolor{ecolite}{RGB}{250,247,242}

\titleformat{\section}{\Large\bfseries\color{ecodark}}{\thesection}{0.7em}{}
\titleformat{\subsection}{\large\bfseries\color{ecoaccent}}{}{0em}{}
\setlist[itemize]{leftmargin=1.4em,itemsep=0.25em,topsep=0.4em}
\setlist[enumerate]{leftmargin=1.6em,itemsep=0.25em,topsep=0.4em}

\newmdenv[
  linecolor=ecowood,linewidth=2pt,topline=false,bottomline=false,rightline=false,
  backgroundcolor=ecolite,leftmargin=0pt,rightmargin=0pt,
  innerleftmargin=1em,innertopmargin=0.7em,innerbottommargin=0.7em,skipabove=1em,skipbelow=1em
]{callout}

\pagestyle{fancy}
\fancyhf{}
\fancyhead[L]{\footnotesize\color{ecodark}EcoWoods Hardwood Flooring \textbullet{} Toronto}
\fancyhead[R]{\footnotesize\color{ecodark}Climate Mastery v1.0}
\fancyfoot[C]{\footnotesize\thepage}
\renewcommand{\headrulewidth}{0.4pt}

\begin{document}
\begin{titlepage}
\vspace*{3cm}
\begin{flushleft}
{\Huge\bfseries\color{ecodark}Climate Mastery}\\[0.6em]
{\LARGE\color{ecoaccent}Why hardwood succeeds or fails in Toronto}\\[2em]
{\large The physical science, moisture protocol and non-negotiable steps that decide whether a hardwood floor becomes a permanent asset or a permanent liability in Toronto's climate.}\\[2.5em]
\rule{\textwidth}{0.6pt}\\[0.8em]
\textbf{Version} 1.0 \quad\textbullet\quad \textbf{Published} 2026-08-01 \quad\textbullet\quad \textbf{Reading time} 9 minutes\\[0.4em]
\textbf{Audience} Homeowners, designers, general contractors, property managers\\[0.4em]
\textbf{Topics} Moisture · Acclimation · Substrate · Installation method · Failure modes\\[1.2em]
\small\textit{The current edition of this paper is published as HTML at }\texttt{ecowoods.ca/papers/toronto-hardwood-climate-moisture-protocol}\textit{. Where the two disagree, the page is correct.}
\end{flushleft}
\end{titlepage}

\tableofcontents
\newpage

\section*{Summary}
\addcontentsline{toc}{section}{Summary}
Toronto's indoor relative humidity swings from the high teens in winter to above sixty percent in summer --- one of the most aggressive annual ranges of any major North American residential market. This paper sets out what that does to wood, the measurements that have to happen before a single board is laid, and the failure modes that appear when any step is skipped.

\section{Toronto's climate reality}
Wood is hygroscopic and anisotropic: it exchanges moisture with the surrounding air continuously, and it moves. Expansion and contraction occur primarily across the grain --- length change is minimal, width change is significant and cumulative. Every board arrives with its own moisture history.

The question is never whether a floor will move. It is whether the floor was specified, tested, acclimated and installed to survive decades of that movement.


\begin{center}
\footnotesize
\begin{longtable}{p{8.00cm}p{8.00cm}}
\caption*{\small\itshape Indoor relative humidity, Toronto residential}\\
\toprule
\textbf{Condition} & \textbf{Relative humidity} \\
\midrule
\endhead
Winter indoor low & 18--25\% RH \\
Summer indoor high & above 60\% RH \\
Safe operating band for hardwood & 35--55\% RH \\
\bottomrule
\end{longtable}
\end{center}
\begin{callout}
\textbf{\color{ecowood}Core truth}\\[0.3em]
Wood will move. The only question is whether the floor was specified, tested, acclimated and installed to survive decades of that movement.
\end{callout}

\section{Moisture testing --- the first non-negotiable}
A floor installed over a wet subfloor, or with boards that have not equalized, will cup, crown, gap or buckle. There is no warranty language that can override physics.

\begin{enumerate}
  \item Subfloor moisture content must be measured.
  \item Flooring material moisture content must be measured.
  \item Both readings must sit within the manufacturer's and EcoWoods' acceptable delta.
  \item Testing occurs at the free estimate, and again immediately before installation.
\end{enumerate}

\section{Solid versus engineered in the GTA}
Solid hardwood is typically 3/4" (19 mm) with a generational wear layer, highly sensitive to RH swings, and best over plywood with a nail-down installation.

Engineered hardwood pairs a real hardwood wear layer with a 90\textdegree{} cross-ply core. That construction is what gives it dimensional stability over concrete, in condominiums, and above radiant heat.

\begin{callout}
\textbf{\color{ecowood}Position}\\[0.3em]
Engineered is the correct specification for the majority of Toronto projects. We specify what the house can support. We never sell what will fail.
\end{callout}

\section{Correct method matched to substrate}
Installation method is not a preference and not a sales option. It is determined by the substrate, the product construction, and the climate load the floor will face for decades.


\begin{center}
\footnotesize
\begin{longtable}{p{8.00cm}p{8.00cm}}
\toprule
\textbf{Method} & \textbf{When it is correct} \\
\midrule
\endhead
Nail-down & Solid hardwood over plywood \\
Glue-down & Engineered over concrete, or in condominiums \\
Floating & Engineered over radiant, or where acoustic separation is required \\
\bottomrule
\end{longtable}
\end{center}
\section{The EcoWoods non-negotiable protocol}
These steps prevent physical failure, and they are what make a lifetime workmanship warranty something a company can offer and still sleep at night.

\begin{enumerate}
  \item Moisture testing of subfloor and material, documented.
  \item Minimum 72-hour acclimation in the actual conditioned space.
  \item Correct installation method matched to substrate.
  \item Proper expansion gaps at all fixed objects and walls.
  \item HEPA dust containment throughout the process.
\end{enumerate}

\section{What happens when a step is skipped}
These are not aesthetic issues. They are permanent, visible records of process failure that no amount of later refinishing can fully erase.

\begin{itemize}
  \item Cupping --- edges higher than the centre of the board
  \item Crowning --- centre higher than the edges
  \item Seasonal gapping
  \item Buckling and tenting
  \item Edge peaking and finish failure
\end{itemize}

\section{What an informed homeowner should demand}
Any company that cannot or will not provide these is optimizing for speed and lowest bid, not for decades of performance.

\begin{enumerate}
  \item Written moisture readings of both subfloor and material before any deposit.
  \item Explicit confirmation of minimum 72-hour acclimation in the actual space.
  \item A fixed price in writing, with no open-ended "unforeseen conditions" language.
  \item Confirmation that salaried artisans, not subcontractors, will perform the work.
  \item Lifetime workmanship warranty language in the contract.
  \item Willingness to refuse the job if conditions are wrong.
\end{enumerate}

\end{document}
